% Figure: CRT bijection Z/15 -> (Z/3)x(Z/5) as a filled grid.
% Source for Figure~\ref{fig:crtgrid}. TikZ only; regenerable by compiling main.tex.
\begin{tikzpicture}[
  cell/.style={draw,minimum size=7.5mm,font=\tiny},
  lab/.style={font=\footnotesize\sffamily,text=slate,align=center},
  hd/.style={font=\footnotesize\sffamily,text=slate}
]
% x mod 15 in cell (row = x mod 3, col = x mod 5)
\foreach \r in {0,1,2} {\foreach \c in {0,1,2,3,4} {
  \pgfmathsetmacro{\v}{int(mod(mod(\c-\r,5)*3*2 - mod(\c-\r,5)*3*0 + \r*10 + \c*0, 15))}
}}
% fill explicitly: x with x mod 3 = r, x mod 5 = c
\foreach \x in {0,...,14} {
  \pgfmathsetmacro{\r}{int(mod(\x,3))}
  \pgfmathsetmacro{\c}{int(mod(\x,5))}
  \node[cell] at (\c*0.78,-\r*0.78) {\x};
}
\foreach \c in {0,1,2,3,4} \node[hd] at (\c*0.78,0.55) {\c};
\node[hd] at (3.2,0.55) {$\leftarrow \bmod\ 5$};
\foreach \r in {0,1,2} \node[hd] at (-0.55,-\r*0.78) {\r};
\node[hd,rotate=90] at (-1.05,-0.78) {$\bmod\ 3$};
\node[lab] at (1.55,-2.9) {Every cell holds exactly one $x \in \{0,\dots,14\}$:\\
CRT says the map $x \mapsto (x \bmod 3,\ x \bmod 5)$ is a bijection.};
\end{tikzpicture}
